Research on reproducing the canonical stylized facts of financial markets divides into bottom-up agent-based models, which derive return dynamics from simulated interactions of heterogeneous agents \cite{lux_marchesi_1999, cont_bouchaud_2000, chen_herding_2015, farmer_joshi_2002, lebaron_agent_2006, arthur_complexity_2021, farmer_quantitative_2025}, and top-down data-driven models, which fit statistical structures directly to observed return series. While agent-based models produce stylized facts as emergent properties, they can be challenging to calibrate to specific assets and computationally demanding as practical generators. The data-driven tradition is the focus of the present review; the subsections below trace the evolution from parametric return models through latent-state and deep generative approaches to the multi-asset extensions that motivate this study.

\subsection{Stylized facts and the limitations of parametric models}
The Black-Scholes-Merton option pricing model \cite{black_scholes_1973} and its underlying geometric Brownian motion (GBM) assume that log-returns are i.i.d.\ Gaussian: prices move smoothly, volatility is constant, and extreme moves are negligibly rare. These properties carry well-known practical advantages: closed-form option prices, simple portfolio variance calculations, and minimal data requirements. However, the Gaussian diffusion paradigm conflicts with a broad set of empirical regularities documented well before the model's widespread adoption. Mandelbrot \cite{mandelbrot_variation_1963} demonstrated that speculative price changes exhibit heavy tails (\textit{leptokurtosis}) beyond Gaussian predictions. Fama \cite{fama_efficient_1970} reviewed evidence that raw return series display negligible linear autocorrelation, consistent with the efficient markets hypothesis. Yet Mandelbrot also observed that large price changes tend to cluster together regardless of sign, a pattern now called \textit{volatility clustering}. Schwert \cite{schwert_why_1989} and others later confirmed this statistically, showing that the autocorrelation of absolute returns decays slowly over time. Cont \cite{cont_empirical_2001} provides a systematic review confirming these regularities across asset classes, geographic markets, and sampling frequencies. Collectively, heavy tails (leptokurtosis), the absence of return predictability, and volatility clustering define the benchmark against which any model of market dynamics must be evaluated.

Engle \cite{engle_1982} introduced autoregressive conditional heteroscedasticity (ARCH) models to address volatility clustering, and Bollerslev \cite{bollerslev_garch_1986} generalized this to the GARCH(p,q) specification. GARCH models are widely used but have well-documented limitations. Heavy tails must be injected by choosing a fat-tailed noise distribution (e.g., Student-t); the model itself does not produce them. There is no representation of discrete market regimes or sudden jumps. Volatility persistence is controlled by a single decay parameter, so the model cannot distinguish between a calm market that slowly heats up and a sudden crash that lingers \cite{diebold_inoue_2001}. Stochastic volatility models \cite{stein_stock_1991, heston_stochastic_1993} treat volatility as a latent process, but require computationally intensive estimation. An alternative approach introduces Poisson-driven jumps directly into the price process. Merton \cite{merton_jump_1976} added a compound Poisson jump component to GBM, producing a Gaussian-mixture marginal distribution that directly addresses leptokurtosis, and Kou \cite{kou_jump-diffusion_2002} extended this framework with asymmetric jump sizes to capture the empirical pattern that crashes tend to be sharper than rallies. Jump-diffusion models are particularly relevant in equity markets, where macroeconomic announcements, earnings surprises, and geopolitical events generate discrete large-magnitude price movements \cite{au_yeung_jump_2020}. However, in their standard form, jump-diffusion models do not provide a mechanism for volatility persistence: the post-jump volatility immediately returns to its baseline level, absent a separate stochastic volatility component. Hidden Markov models offer a structurally different solution by embedding these dynamics in a latent discrete-state process.

\subsection{Hidden Markov models and regime-switching}
Hamilton \cite{hamilton_1989} introduced the regime-switching HMM to economics, showing that U.S.\ GNP growth is well described by a two-state model with distinct expansion and recession regimes. Since then, HMMs have become a standard tool for latent-state modeling of financial time series, with applications ranging from volatility forecasting \cite{rossi_volatility_2006} to equity trading strategies \cite{nguyen_hidden_2018}. Ryd\'{e}n, Ter\"{a}svirta, and \AA sbrink \cite{ryden_stylized_1998} systematically evaluated HMMs against the stylized facts of daily returns. Gaussian-emission HMMs with few states reproduce heavy tails and the absence of return autocorrelation, but fail to capture volatility clustering: the ACF of absolute returns decays far too quickly relative to empirical data. Bulla and Bulla \cite{bulla_stylized_2006} showed that hidden semi-Markov models (HSMMs), which replace the geometric sojourn-time distribution with a flexible alternative, substantially improve volatility-clustering reproduction by allowing realistic persistence in high-volatility states. However, both standard and semi-Markov approaches rely on EM estimation (Baum-Welch \cite{baum_1970}), which requires iterative forward-backward passes and is sensitive to initialization. The present study sidesteps this limitation by replacing EM with direct frequentist counting of transitions between quantile-defined states, an approach that is computationally trivial and initialization-free. Alongside HMM-based methods, a parallel line of work has explored whether deep generative models can learn stylized facts directly from data.

\subsection{Synthetic data generation and deep generative models}
Synthetic financial time series are increasingly used to augment limited data, stress test portfolios, and benchmark generative models. Assefa et al.\ \cite{assefa_generating_2021} surveyed this landscape, noting that stylized-facts reproduction is the primary quality criterion and that naive generators tend to fail on at least one core stylized fact. Deep generative models, particularly generative adversarial networks (GANs), have attracted interest as a way to learn these properties directly from data, but results have been mixed. Takahashi, Chen, and Tanaka-Ishii \cite{takahashi_2019} found that standard GAN architectures produced return sequences that visually resembled financial data yet failed to reproduce ACF structure, and Kwon and Lee \cite{kwon_can_2024} confirmed that volatility clustering remained elusive unless the architecture was specifically designed for it. Foundation models for Markov jump processes \cite{berghaus_foundation_2025} offer a complementary pretrained-inference perspective, though their application to financial series remains open. The hybrid HMM of the present study takes a different path: it replaces iterative EM with direct frequentist counting (eliminating training instability) and augments the Markov chain with a jump-duration mechanism that addresses the volatility-clustering gap without the architectural complexity of deep generative models. All of the models discussed so far are univariate; extending them to multi-asset portfolios introduces a separate set of challenges.

\subsection{Factor models and multi-asset extensions}
Existing HMM-based work in finance has been predominantly univariate \cite{hamilton_1989, rossi_volatility_2006, nguyen_hidden_2018, kim_nelson_1999}, leaving open the question of how HMM-derived generative models can scale to large cross-sections of assets. One strategy could be to use an HMM-based approach to capture the marginal distribution of a market factor (e.g., SPY) and then reconstruct individual asset paths via a linear factor model. The Single-Index Model (SIM) of Sharpe \cite{sharpe_simplified_1963} decomposes each asset's excess return into a common market factor and an idiosyncratic residual, reducing the parameter count relative to a full multivariate model. Mixture hidden Markov models \cite{dias_mixture_2010} allow multiple assets to share a latent state structure, but joint estimation across hundreds of assets is computationally demanding. On the other hand, copulas offer an alternative by separating marginal distributions from the dependence structure \cite{embrechts_copulas_2002, cherubini_copula_2004}, though estimation grows expensive in high dimensions and the choice of copula family adds model selection complexity. The present study addresses the multi-asset challenge through both approaches: a SIM factor decomposition provides a scalable baseline for a 424-asset universe (the hybrid HMM generates market-factor paths from SPY, and asset-level paths are reconstructed via the linear projection), while copula-based dependence models (Gaussian, Student-t, and C-vine \cite{aas_pair_2009}) preserve each asset's independently fitted HMM marginal and inject cross-asset dependence via rank reordering. The empirical comparison reveals that copulas substantially improve per-asset distributional accuracy relative to the SIM, and that a single Student-t copula is a parsimonious and sufficient specification for the equity pairs examined.
